% !TEX root = ../main.tex
\corrigesChapitre{23}{Effet Doppler}
\corrigesSouples

\begin{blocLivre}{cor-23-1}
\begin{corrige}{1}{La sirène des pompiers}
\demarche{$f_\mathrm{R}=f_\mathrm{E}\,\dfrac{v}{v-v_\mathrm{E}}$ à l'approche, $f_\mathrm{E}\,\dfrac{v}{v+v_\mathrm{E}}$ à l'éloignement, avec $v_\mathrm{E}=\qty{20}{m.s^{-1}}$.}
\begin{questions}
\item À l'approche, les deux notes paraissent \res{\text{plus aiguës}} ; au passage, elles deviennent brusquement \res{\text{plus graves}}.
\item À l'approche : $435\times\dfrac{340}{320}=\res{\qty{462}{Hz}}$ ; à l'éloignement : $435\times\dfrac{340}{360}=\res{\qty{411}{Hz}}$.
\item $488\times\dfrac{340}{320}=\res{\qty{519}{Hz}}$.
\item \res{\text{Non}} : la sirène est immobile par rapport à eux ; ils entendent les fréquences émises.
\end{questions}
\end{corrige}
\end{blocLivre}
\begin{blocLivre}{cor-23-2}
\begin{corrige}{2}{Un émetteur qui s'éloigne}
\demarche{Pendant $T_\mathrm{E}$, le premier maximum avance de $v\,T_\mathrm{E}$ et l'émetteur recule de $v_\mathrm{E}\,T_\mathrm{E}$.}
\begin{questions}
\item L'émetteur parcourt $v_\mathrm{E}\,T_\mathrm{E}$, \res{\text{en s'éloignant}} du récepteur.
\item À la date $T_\mathrm{E}$, le premier maximum est à $v\,T_\mathrm{E}$ devant la position initiale, l'émetteur à $v_\mathrm{E}\,T_\mathrm{E}$ derrière :
$\res{\lambda_\mathrm{R}=(v+v_\mathrm{E})\,T_\mathrm{E}}$.
\item $f_\mathrm{R}=\dfrac{v}{\lambda_\mathrm{R}}=\res{f_\mathrm{E}\,\dfrac{v}{v+v_\mathrm{E}}}$.
\item $\dfrac{v}{v+v_\mathrm{E}}<1$ : $f_\mathrm{R}$ est \res{\text{plus petite}}. Sur le schéma, les fronts émis vers l'arrière sont plus espacés.
\end{questions}
\end{corrige}
\end{blocLivre}
\begin{blocLivre}{cor-23-3}
\begin{corrige}{3}{Le sifflet du train}
\demarche{$v_\mathrm{E}=\dfrac{144}{3{,}6}=\qty{40}{m.s^{-1}}$.}
\begin{questions}
\item $v_\mathrm{E}=\res{\qty{40}{m.s^{-1}}}$.
\item Approche : $500\times\dfrac{340}{300}=\res{\qty{566,7}{Hz}}$ ; éloignement : $500\times\dfrac{340}{380}=\res{\qty{447,4}{Hz}}$.
\item Rapport $\dfrac{566{,}7}{447{,}4}=\res{1{,}27}$ : un écart bien \res{\text{perceptible}}, voisin d'un intervalle de tierce en musique.
\item \res{\text{Non}} : le sifflet est immobile par rapport au conducteur, qui entend \qty{500}{Hz}.
\end{questions}
\end{corrige}
\end{blocLivre}
\begin{blocLivre}{cor-23-4}
\begin{corrige}{4}{Quand l'approximation est-elle valable ?}
\demarche{$\Delta f_\text{exact}=f_\mathrm{E}\,\dfrac{v_\mathrm{E}}{v-v_\mathrm{E}}$ ; $\Delta f_\text{approché}=f_\mathrm{E}\,\dfrac{v_\mathrm{E}}{v}$.}
\begin{questions}
\item \res{\qty{30,3}{Hz}} ; \res{\qty{172}{Hz}} ; \res{\qty{416,7}{Hz}}.
\item \res{\qty{29,4}{Hz}} ; \res{\qty{147}{Hz}} ; \res{\qty{294}{Hz}}.
\item À \qty{10}{m.s^{-1}}, l'écart relatif vaut \res{\qty{2,9}{\percent}} : l'approximation n'est valable que si $v_\mathrm{E}$ est
\res{\text{très petite}} devant $v$ (à quelques pour cent près).
\item Pour la lumière, $v=c=\qty{3e8}{m.s^{-1}}$ : $v_\mathrm{E}/c$ est minuscule pour tout véhicule, et l'approximation est
\res{\text{excellente}}.
\end{questions}
\end{corrige}
\end{blocLivre}
\begin{blocLivre}{cor-23-5}
\begin{corrige}{5}{Une étoile qui s'éloigne}
\demarche{$v_\mathrm{E}\approx c\,\dfrac{\lambda_\mathrm{R}-\lambda_\mathrm{E}}{\lambda_\mathrm{E}}$.}
\begin{questions}
\item $\lambda_\mathrm{R}>\lambda_\mathrm{E}$ : décalage \res{\text{vers le rouge}}, l'étoile \res{\text{s'éloigne}}.
\item $v_\mathrm{E}=\num{3.00e8}\times\dfrac{0{,}22}{656{,}28}=\qty{1,01e5}{m.s^{-1}}$, soit \res{\qty{101}{km.s^{-1}}}.
\item Le décalage relatif n'est que de \num{3e-4} : une erreur de \qty{0,01}{nm} change la vitesse de près de \qty{5}{km.s^{-1}}.
\item Seule la composante de la vitesse \res{\text{dans la direction de la Terre}} (vitesse radiale).
\end{questions}
\end{corrige}
\end{blocLivre}
\begin{blocLivre}{cor-23-6}
\begin{corrige}{6}{Le radar routier}
\demarche{$v_\mathrm{E}=\dfrac{c\,\Delta f}{2\,f\cos\alpha}$.}
\begin{questions}
\item L'onde subit un premier décalage en atteignant la voiture, puis un second en revenant : d'où le facteur
\res{2}. Seule la composante de la vitesse \res{\text{dans l'axe du faisceau}}, $v_\mathrm{E}\cos\alpha$, intervient.
\item $v_\mathrm{E}=\dfrac{\num{3.00e8}\times3950}{2\times\num{24.125e9}\times\cos25^\circ}=\res{\qty{27,1}{m.s^{-1}}}$, soit \res{\qty{97,6}{km.h^{-1}}}.
\item Vitesse retenue : \qty{92,6}{km.h^{-1}} $>$ \qty{90}{km.h^{-1}} : le conducteur est \res{\text{en infraction}}.
\item $\Delta f=\dfrac{2\times\num{24.125e9}\times25\times\cos25^\circ}{\num{3.00e8}}=\res{\qty{3,64}{kHz}}$.
\end{questions}
\end{corrige}
\end{blocLivre}
\begin{blocLivre}{cor-23-7}
\begin{corrige}{7}{La vitesse d'une voiture au bord de la route}
\demarche{Les deux expressions de $f_\mathrm{R}$ font intervenir la même $f_\mathrm{E}$ : on l'élimine.}
\begin{questions}
\item $f_\mathrm{E}=f_\text{app}\,\dfrac{v-v_\mathrm{E}}{v}=f_\text{éloi}\,\dfrac{v+v_\mathrm{E}}{v}$, d'où la relation ; en développant,
$v_\mathrm{E}\,(f_\text{app}+f_\text{éloi})=v\,(f_\text{app}-f_\text{éloi})$ : c'est la ligne de \code{vE}.
\item C'est la relation $f_\mathrm{E}=f_\text{app}\,\dfrac{v-v_\mathrm{E}}{v}$.
\item On remplace \code{971.4} et \code{838.4} par \code{1062} et \code{878} : $v_\mathrm{E}=\res{\qty{32,2}{m.s^{-1}}}$, soit
\res{\qty{116}{km.h^{-1}}}, et $f_\mathrm{E}=\res{\qty{961}{Hz}}$.
\item Au moment du passage, la voiture ne se déplace plus dans la direction de l'élève : la fréquence varie
rapidement et \res{\text{ne vérifie pas}} les relations à une dimension.
\end{questions}
\end{corrige}
\end{blocLivre}
\begin{blocLivre}{cor-23-8}
\begin{corrige}{8}{L'échographie Doppler}
\demarche{$v=\dfrac{c_\mathrm{t}\,\Delta f}{2\,f\cos\theta}$.}
\begin{questions}
\item $v=\dfrac{1540\times1300}{2\times\num{4.0e6}\times\cos60^\circ}=\res{\qty{0,50}{m.s^{-1}}}$.
\item \qty{1,3}{kHz} est \res{\text{audible}} : le médecin entend directement l'écoulement et repère les anomalies au son.
\item $\Delta f=\dfrac{2\times\num{4.0e6}\times1{,}0\times0{,}5}{1540}=\res{\qty{2,6}{kHz}}$.
\item Pour $\theta=90^\circ$, $\cos\theta=0$ : le décalage est \res{\text{nul}}, quelle que soit la vitesse du sang.
\end{questions}
\end{corrige}
\end{blocLivre}
\begin{blocLivre}{cor-23-9}
\begin{corrige}{9}{Vers le mur du son}
\demarche{$\lambda_\mathrm{R}=\dfrac{v-v_\mathrm{E}}{f_\mathrm{E}}$ devant l'avion.}
\begin{questions}
\item À l'arrêt : $\lambda=\dfrac{340}{1000}=\res{\qty{34}{cm}}$ ; à $0{,}9\,v$ : $\lambda_\mathrm{R}=\dfrac{0{,}1\times340}{1000}=\res{\qty{3,4}{cm}}$.
\item $f_\mathrm{R}=\dfrac{v}{\lambda_\mathrm{R}}=\res{\qty{10}{kHz}}$.
\item $\lambda_\mathrm{R}$ \res{\text{tend vers 0}} : les fronts d'onde s'accumulent devant l'avion et se superposent.
\item Au-delà de la vitesse du son, ces fronts accumulés forment une onde de choc, une surpression brutale qui balaie le
sol derrière l'avion : c'est le \res{\text{bang}}.
\end{questions}
\end{corrige}
\end{blocLivre}
\begin{blocLivre}{cor-23-10}
\begin{corrige}{10}{Une étoile double}
\demarche{L'écart maximal à la raie de laboratoire donne la vitesse radiale maximale.}
\begin{questions}
\item L'étoile tourne autour du centre de masse du système : elle se rapproche et s'éloigne \res{\text{alternativement}}.
\item Elle se rapproche quand la raie est à \qty{656,10}{nm} (décalage vers le bleu), s'éloigne quand elle est à
\qty{656,46}{nm} (vers le rouge).
\item $v_\mathrm{max}=\num{3.00e8}\times\dfrac{0{,}18}{656{,}28}=\res{\qty{82}{km.s^{-1}}}$.
\item La vitesse est alors constamment \res{\text{perpendiculaire}} à la direction de la Terre : aucune composante radiale.
\end{questions}
\end{corrige}
\end{blocLivre}
\begin{blocLivre}{cor-23-11}
\begin{corrige}{11}{Le son d'une moto}
\demarche{Comme dans le programme de l'essentiel : $v_\mathrm{E}=v\,\dfrac{f_\text{app}-f_\text{éloi}}{f_\text{app}+f_\text{éloi}}$.}
\begin{questions}
\item À l'approche : \res{\qty{330}{Hz}} ; à l'éloignement : \res{\qty{270}{Hz}}.
\item $f_\text{app}(v-v_\mathrm{E})=f_\text{éloi}(v+v_\mathrm{E})$ donne $v_\mathrm{E}=340\times\dfrac{60}{600}=\res{\qty{34}{m.s^{-1}}}$, soit
\res{\qty{122}{km.h^{-1}}}.
\item $f_\mathrm{E}=330\times\dfrac{340-34}{340}=\res{\qty{297}{Hz}}$.
\item Vers $t=\res{\qty{3}{s}}$, quand la fréquence chute brutalement et passe par la fréquence émise.
\end{questions}
\end{corrige}
\end{blocLivre}
\begin{blocLivre}{cor-23-12}
\begin{corrige}{12}{L'expansion de l'Univers}
\demarche{$v=c\,\dfrac{\lambda_\mathrm{R}-\lambda_\mathrm{E}}{\lambda_\mathrm{E}}$, puis rapport $v/d$.}
\begin{questions}
\item A : \res{\qty{1417}{km.s^{-1}}} ; B : \res{\qty{3474}{km.s^{-1}}} ; C : \res{\qty{7039}{km.s^{-1}}}.
\item $v/d$ vaut 70,9 ; 69,5 et \qty{70,4}{km.s^{-1}} par mégaparsec : la vitesse est bien \res{\text{proportionnelle}} à la distance,
avec $\res{H_0\approx\qty{70}{km.s^{-1}}}$ par mégaparsec.
\item $\dfrac{1}{H_0}=\dfrac{\num{3.086e19}}{70}=\qty{4,4e17}{s}$, soit environ \res{\text{14 milliards d'années}}.
\item $\Delta\lambda/\lambda\approx0{,}11$ : la vitesse n'est plus très petite devant $c$, l'approximation devient
\res{\text{insuffisante}}.
\end{questions}
\end{corrige}
\end{blocLivre}
\begin{blocLivre}{cor-23-13}
\begin{corrige}{13}{À la recherche des exoplanètes}
\demarche{Décalage relatif $\dfrac{\Delta\lambda}{\lambda}=\dfrac{v}{c}$.}
\begin{questions}
\item L'étoile tourne autour du centre de masse commun : ses raies oscillent \res{\text{périodiquement}} autour de leur position,
avec la période de l'orbite de la planète.
\item $\dfrac{56}{\num{3.00e8}}=\res{\num{1.9e-7}}$, soit $\Delta\lambda=\res{\qty{9,3e-5}{nm}}$ pour une raie de \qty{500}{nm}.
\item Le décalage relatif dépasse d'un facteur 20 la précision de $10^{-8}$ : il est \res{\text{détectable}}, mais il correspond
à un déplacement de raie d'un dix-millième de nanomètre, ce qui exige un spectrographe d'une stabilité extrême.
\item Le décalage serait 4,5 fois plus faible, avec une période de 12~ans : \res{\text{plus difficile}} à détecter.
\end{questions}
\end{corrige}
\end{blocLivre}
\begin{blocLivre}{cor-23-14}
\begin{corrige}{14}{Écouter un satellite}
\demarche{$\Delta f=f\,\dfrac{v_\mathrm{r}}{c}$, avec la composante radiale $v_\mathrm{r}$ de la vitesse.}
\begin{questions}
\item $\Delta f_\mathrm{max}=\num{437e6}\times\dfrac{7000}{\num{3.00e8}}=\res{\qty{10,2}{kHz}}$ ; l'excursion totale vaut \res{\qty{20,4}{kHz}}.
\item Le décalage dépasse la bande passante du récepteur : sans ajustement, le signal \res{\text{sortirait de la bande}} écoutée.
\item À \qty{2,4}{GHz}, $\Delta f_\mathrm{max}=\qty{56}{kHz}$ : le problème est \res{\text{plus gênant}}, le décalage étant proportionnel à la
fréquence.
\item Le satellite se déplace alors perpendiculairement à la direction de la station : la vitesse radiale est nulle et le
décalage \res{\text{aussi}}.
\end{questions}
\end{corrige}
\end{blocLivre}
\begin{blocLivre}{cor-23-15}
\begin{corrige}{15}{Le radar météorologique}
\demarche{$v_\mathrm{r}=\dfrac{c\,\Delta f}{2\,f}$, positive si les gouttes se rapprochent du radar.}
\begin{questions}
\item $v_\mathrm{r}=\dfrac{\num{3.00e8}\times560}{2\times\num{5.6e9}}=\res{\qty{15}{m.s^{-1}}}$, soit \res{\qty{54}{km.h^{-1}}}.
\item Au nord, les gouttes se rapprochent du radar : elles se dirigent vers le sud. Au sud, elles s'éloignent : elles se
dirigent aussi vers le sud. Le vent souffle \res{\text{du nord vers le sud}}, à \qty{54}{km.h^{-1}}.
\item $\Delta f=\dfrac{2\times\num{5.6e9}\times8}{\num{3.00e8}}=\res{\qty{299}{Hz}}$.
\item La vitesse radiale est alors \res{\text{nulle}} : le radar ne mesure que la composante du vent dans la direction de visée.
\end{questions}
\end{corrige}
\end{blocLivre}
